\documentclass[mode=handout, palette=emerald, lang=french, fontsize=11pt]{modernclassnotes}

\course{Mécanique Quantique et Physique Statistique}
\coursecode{PHYS 301}
\professor{Prof. Claire Dubois}
\institution{Sorbonne Université}
\term{Semestre d'Automne 2026}
\lecturedate{12 Novembre 2026}
\title{Équation de Schrödinger et Puits de Potentiel}
\docsubtitle{Fiche de synthèse – Cours 3}

\begin{document}

\maketitle

\section{Introduction et Principes Fondamentaux}

L'état quantique d'une particule est décrit par sa fonction d'onde $\Psi(x,t)$, solution de l'équation de Schrödinger dépendante du temps.

\begin{definition}[Équation de Schrödinger Indépendante du Temps]
Pour un système conservatif d'énergie $E$, l'équation aux valeurs propres s'écrit :
\begin{equation*}
  \hat{H} \psi(x) = -\frac{\hbar^2}{2m} \frac{d^2\psi(x)}{dx^2} + V(x)\psi(x) = E \psi(x)
\end{equation*}
où $\hat{H}$ est l'opérateur hamiltonien.
\end{definition}

\begin{theorem}[Quantification de l'Énergie dans un Puits Infini]
Pour une particule confinée dans un puits potentiel unidimensionnel de largeur $L$ avec $V(x)=0$ sur $[0,L]$ et $V(x)=\infty$ ailleurs, les niveaux d'énergie admissibles sont quantifiés :
\begin{equation*}
  E_n = \frac{n^2 \pi^2 \hbar^2}{2m L^2}, \quad n \in \{1, 2, 3, \dots\}
\end{equation*}
\end{theorem}

\begin{remark}
Notez que l'état fondamental ($n=1$) possède une énergie non nulle $E_1 > 0$, appelée \textbf{énergie du point zéro}, imposée par le principe d'incertitude d'Heisenberg $\Delta x \cdot \Delta p \ge \frac{\hbar}{2}$.
\end{remark}

\begin{tip}[Méthode de Résolution]
Toujours vérifier la continuité de la fonction d'onde $\psi(x)$ et de sa dérivée première $\psi'(x)$ aux limites des régions de potentiel fini.
\end{tip}

\begin{takeaways}
\begin{itemize}
  \item La fonction d'onde doit être continue, dérivable et de carré intégrable ($\int |\psi|^2 dx = 1$).
  \item Le confinement spatial entraîne la quantification automatique des états énergétiques.
  \item L'effet tunnel apparaît dès que le potentiel $V(x) > E$ sur une barrière de largeur finie.
\end{itemize}
\end{takeaways}

\end{document}
